\chapter{Conclusioni}
\section{Discussione critica dei risultati ottenuti}
\begin{itemize}
	\item sviluppo del codice firmware, durata media e massima di un loop
	\item taratura dei parametri e conoscenza dei limiti del controllore PID
	\item miglioramenti del pid rispetto al controllore ad isteresi
\end{itemize}

\section{Possibili sviluppi futuri}
\begin{itemize}
	\item implementazione di un sistema di controllo dell'umidità (o controllo multivariabile)
	\item valutare se implementare il pid per farlo lavorare su interi e non float, stando attenti agli overflow
	\item utilizzare mosfet al posto dei relè e aggiungere un sistema di smorzamento dell'output alla riattivazione
	del pid
	\item modellare il sistema per implementare un controllore dal punto di vista teorico con simulazioni e poi
	testare i parametri ottimali sul prototipo reale
	\item integrare un'azione di feedforward (considerando l'inerzia termica) e attivare il pid solo a regime
	\item campo di prova per implementare e testare altri tipi di controllore (ad esempio fuzzy, predittivo, ecc.)
\end{itemize}
